\ifdefined\gkver\else\def\gkver{student}\fi
\documentclass[\gkver]{gaokaozhenti}
\begin{document}

\chapter{2023年辽宁}

\begin{enumerate}
\item
%% number: 1
%% paperName: 2023年辽宁高考第 1 题 4 分
%% typeId: 1
%% type: 单选题
%% chapter: 曲线运动
%% point: 曲线运动条件
%% method: 无
%% score: 4
%% degree: 850
%% duplicateId: 0
%% body:
某同学在练习投篮，篮球在空中的运动轨迹如图中虚线所示，篮球所受合力 $F$ 的示意图可能正确的是\xzanswer{A}

\onepicture[align=right, w=3.5cm]{figs/01a.png}

\onepicture[w=13cm]{\includesvg{figs/01b}}
%% answer: A
%% memo:
\memoanswer{
	篮球做曲线运动，所受合力指向运动轨迹的凹侧。

	故选 A。
}

\item
%% number: 2
%% paperName: 2023年辽宁高考第 2 题 4 分
%% typeId: 1
%% type: 单选题
%% chapter: 运动和力
%% point: 力学单位制
%% method: 无
%% score: 4
%% degree: 750
%% duplicateId: 0
%% body:
安培通过实验研究，发现了电流之间相互作用力的规律。若两段长度分别为 $\Delta l_{1}$ 和 $\Delta l_{2}$、电流大小分别为 $I_{1}$ 和 $I_{2}$ 的平行直导线间距为 $r$ 时，相互作用力的大小可以表示为 $\Delta F = k \frac{{{I_1}{I_2}\Delta {l_1}\Delta {l_2}}}{{{r^2}}}$。比例系数 $k$ 的单位是\xzanswer{B}

\fourchoices[answer=B]
{$\Ukg\cdot\Um/(\Usq\cdot\UA)$}
{$\Ukg\cdot\Um/(\Usq\cdot\UA^{2})$}
{$\Ukg\cdot\Umq/(\Us^{3}\cdot\UA)$}
{$\Ukg\cdot\Umq/(\Us^{3}\cdot\UA^{3})$}
%% answer: B
%% memo:
\memoanswer{
	根据题干公式 $\Delta F = k \frac{I_{1}I_{2}\Delta l_{1}\Delta l_{2}}{r^{2}}$ 整理可得 $k = \frac{\Delta F r^{2}}{I_{1}I_{2}\Delta l_{1}\Delta l_{2}}$

	代入相应物理量单位可得比例系数 $k$ 的单位为 $\frac{\UN}{\UA^{2}} = \frac{\Ukg\cdot\Um/\Usq}{\UA^{2}} = \Ukg\cdot\Um/(\Usq\cdot\UA^{2})$。

	故选 B。
}

\item
%% number: 3
%% paperName: 2023年辽宁高考第 3 题 4 分
%% typeId: 1
%% type: 单选题
%% chapter: 直线运动
%% point: 运动图象
%% method: 无
%% score: 4
%% degree: 550
%% duplicateId: 0
%% body:
如图（a），从高处 M 点到地面 N 点有 Ⅰ、Ⅱ 两条光滑轨道。两相同小物块甲、乙同时从 M 点由静止释放，沿不同轨道滑到 N 点，其速率 $v$ 与时间 $t$ 的关系如图（b）所示。由图可知，两物块在离开 M 点后、到达 N 点前的下滑过程中\xzanswer{B}
\onepicture[w=9cm]{figs/03.png}
\fourchoices[answer=B]
{甲沿 I 下滑且同一时刻甲的动能比乙的大}
{甲沿 Ⅱ 下滑且同一时刻甲的动能比乙的小}
{乙沿 I 下滑且乙的重力功率一直不变}
{乙沿 Ⅱ 下滑且乙的重力功率一直增大}
%% answer: B
%% memo:
\memoanswer{
	AB．由图乙可知，甲下滑过程中，甲做匀加速直线运动，则甲沿 Ⅱ 下滑，乙做加速度逐渐减小的加速运动，乙沿 I 下滑，任意时刻甲的速度都小于乙的速度，可知同一时刻甲的动能比乙的小，A 错误，B 正确；

	CD．乙沿 I 下滑，开始时乙速度为 0，到 N 点时乙竖直方向速度为零，根据瞬时功率公式 $P = mgv_{y}$ 可知重力瞬时功率先增大后减小，CD 错误。

	故选 B。
}

\item
%% number: 4
%% paperName: 2023年辽宁高考第 4 题 4 分
%% typeId: 1
%% type: 单选题
%% chapter: 电磁感应
%% point: 切割磁感线电动势
%% method: 无
%% score: 4
%% degree: 650
%% duplicateId: 0
%% body:
如图，空间中存在水平向右的匀强磁场，一导体棒绕固定的竖直轴 OP 在磁场中匀速转动，且始终平行于 OP。导体棒两端的电势差 $u$ 随时间 $t$ 变化的图像可能正确的是\xzanswer{C}

\onepicture[w=9cm]{figs/04a.png}

\onepicture[w=13cm]{\includesvg{figs/04b}}
%% answer: C
%% memo:
\memoanswer{
	如图所示，导体棒匀速转动，设速度为 $v$，设导体棒从 A 到 B 过程，棒穿过的角度为 $\theta$，则导体棒垂直磁感线方向的分速度为 $v=v_{\perp}\cos\theta$，可知导体棒垂直磁感线的分速度为余弦变化，根据左手定则可知，导体棒经过 B 点和 B 点关于 P 点的对称点时，电流方向发生变化，根据 $u=BLv_{\perp}$ 可知导体棒两端的电势差 $u$ 随时间 $t$ 变化的图像为余弦图像。

	故选 C。
}

\item
%% number: 5
%% paperName: 2023年辽宁高考第 5 题 4 分
%% typeId: 1
%% type: 单选题
%% chapter: 气体性质
%% point: 气体图象
%% method: 无
%% score: 4
%% degree: 550
%% duplicateId: 0
%% body:
“空气充电宝”是一种通过压缩空气实现储能的装置，可在用电低谷时储存能量、用电高峰时释放能量。“空气充电宝”某个工作过程中，一定质量理想气体的 $p - T$ 图像如图所示。该过程对应的 $p - V$ 图像可能是\xzanswer{B}

\onepicture[w=3.5cm]{figs/05a.png}

\onepicture[w=13cm]{\includesvg{figs/05b}}
%% answer: B
%% memo:
\memoanswer{
	从 a 到 b，气体压强不变，温度升高，则体积变大；从 b 到 c，气体压强减小，温度降低，因 c 点与原点连线的斜率小于 b 点与原点连线的斜率，c 态的体积大于 b 态体积。

	故选 B。
}

\item
%% number: 6
%% paperName: 2023年辽宁高考第 6 题 4 分
%% typeId: 1
%% type: 单选题
%% chapter: 原子和原子核
%% point: 玻尔模型
%% method: 无
%% score: 4
%% degree: 650
%% duplicateId: 0
%% body:
原子处于磁场中，某些能级会发生劈裂。某种原子能级劈裂前后的部分能级图如图所示，相应能级跃迁放出的光子分别设为①②③④。若用①照射某金属表面时能发生光电效应，且逸出光电子的最大初动能为 $E_{k}$，则\xzanswer{A}
\onepicture[w=6cm]{figs/06.png}
\fourchoices[answer=A]
{①和③的能量相等}
{②的频率大于④的频率}
{用②照射该金属一定能发生光电效应}
{用④照射该金属逸出光电子的最大初动能小于 $E_{k}$}
%% answer: A
%% memo:
\memoanswer{
	A．由图可知①和③对应的跃迁能级差相同，可知①和③的能量相等，选项 A 正确；

	B．因②对应的能级差小于④对应的能级差，可知②的能量小于④的能量，根据 $E = h\nu$ 可知②的频率小于④的频率，选项 B 错误；

	C．因②对应的能级差小于①对应的能级差，可知②的能量小于①，②的频率小于①，则若用①照射某金属表面时能发生光电效应，用②照射该金属不一定能发生光电效应，选项 C 错误；

	D．因④对应的能级差大于①对应的能级差，可知④的能量大于①，即④的频率大于①，因用①照射某金属表面时能逸出光电子的最大初动能为 $E_{k}$，根据 $E_{km} = h\nu - W_{\text{逸出功}}$，则用④照射该金属逸出光电子的最大初动能大于 $E_{k}$，选项 D 错误。

	故选 A。
}

\item
%% number: 7
%% paperName: 2023年辽宁高考第 7 题 4 分
%% typeId: 1
%% type: 单选题
%% chapter: 曲线运动
%% point: 天体运动
%% method: 无
%% score: 4
%% degree: 550
%% duplicateId: 0
%% body:
在地球上观察，月球和太阳的角直径（直径对应的张角）近似相等，如图所示。若月球绕地球运动的周期为 $T_{1}$，地球绕太阳运动的周期为 $T_{2}$，地球半径是月球半径的 $k$ 倍，则地球与太阳的平均密度之比约为\xzanswer{D}
\onepicture[w=9cm]{figs/07.png}
\fourchoices[answer=D]
{$k^{3}\left(\frac{T_{1}}{T_{2}}\right)^{2}$}
{$k^{3}\left(\frac{T_{2}}{T_{1}}\right)^{2}$}
{$\frac{1}{k^{3}}\left(\frac{T_{1}}{T_{2}}\right)^{2}$}
{$\frac{1}{k^{3}}\left(\frac{T_{2}}{T_{1}}\right)^{2}$}
%% answer: D
%% memo:
\memoanswer{
	设月球绕地球运动的轨道半径为 $r_{1}$，地球绕太阳运动的轨道半径为 $r_{2}$，根据 $G\frac{Mm}{r^{2}}=m\frac{4\pi^{2}}{T^{2}}r$

	可得

	$G\frac{m_{\text{地}}m_{\text{月}}}{r_{1}^{2}}=m_{\text{月}}\frac{4\pi^{2}}{T_{1}^{2}}r_{1}$

	$G\frac{m_{\text{地}}m_{\text{日}}}{r_{2}^{2}}=m_{\text{地}}\frac{4\pi^{2}}{T_{2}^{2}}r_{2}$

	其中

	$\frac{r_{1}}{r_{2}}=\frac{R_{\text{月}}}{R_{\text{日}}}=\frac{R_{\text{地}}}{kR_{\text{日}}}$

	$\rho=\frac{m}{\frac{4}{3}\pi R^{3}}$

	联立可得 $\frac{\rho_{\text{地}}}{\rho_{\text{日}}} = \frac{1}{k^{3}} \left( \frac{T_{2}}{T_{1}} \right)^{2}$

	故选 D。
}

\item
%% number: 8
%% paperName: 2023年辽宁高考第 8 题 6 分
%% typeId: 2
%% type: 多选题
%% chapter: 机械波
%% point: 波的叠加
%% method: 无
%% score: 6
%% degree: 750
%% duplicateId: 0
%% body:
“球鼻艏”是位于远洋轮船船头水面下方的装置，当轮船以设计的标准速度航行时，球鼻艏推起的波与船首推起的波如图所示，两列波的叠加可以大幅度减小水对轮船的阻力。下列现象的物理原理与之相同的是\xzanswer{BD}
\onepicture[w=7cm]{figs/08.png}
\fourchoices[answer=BD]
{插入水中的筷子、看起来折断了}
{阳光下的肥皂膜，呈现彩色条纹}
{驶近站台的火车，汽笛音调变高}
{振动音叉的周围，声音忽高忽低}
%% answer: BD
%% memo:
\memoanswer{
	该现象属于波的叠加原理；插入水中的筷子看起来折断了是光的折射造成的，与该问题的物理原理不相符；阳光下的肥皂膜呈现彩色条纹，是由于光从薄膜上下表面的反射光叠加造成的干涉现象，与该问题的物理原理相符；驶近站台的火车汽笛音调变高是多普勒现象造成的，与该问题的物理原理不相符；振动音叉的周围声音忽高忽低，是声音的叠加造成的干涉现象，与该问题的物理原理相符。

	故选 BD。
}

\item
%% number: 9
%% paperName: 2023年辽宁高考第 9 题 6 分
%% typeId: 1
%% type: 单选题
%% chapter: 电场
%% point: 等势面
%% method: 无
%% score: 6
%% degree: 550
%% duplicateId: 0
%% body:
图（a）为金属四极杆带电粒子质量分析器的局部结构示意图，图（b）为四极杆内垂直于 $x$ 轴的任意截面内的等势面分布图，相邻两等势面间电势差相等，则\xzanswer{CD}
\onepicture[w=10cm]{figs/09.png}
\fourchoices[answer=CD]
{P 点电势比 M 点的低}
{P 点电场强度大小比 M 点的大}
{M 点电场强度方向沿 $z$ 轴正方向}
{沿 $x$ 轴运动的带电粒子，电势能不变}
%% answer: CD
%% memo:
\memoanswer{
	A．因 P 点所在的等势面高于 M 点所在的等势面，可知 P 点电势比 M 点的高，选项 A 错误；

	B．因 M 点所在的等差等势面密集，则 M 点场强较 P 点大，即 P 点电场强度大小比 M 点的小，选项 B 错误；

	C．场强方向垂直等势面，且沿电场线方向电势逐渐降低，可知 M 点电场强度方向沿 $z$ 轴正方向，选项 C 正确；

	D．因 $x$ 轴上各点电势相等，则沿 $x$ 轴运动的带电粒子，电势能不变，选项 D 正确。

	故选 CD。
}

\item
%% number: 10
%% paperName: 2023年辽宁高考第 10 题 6 分
%% typeId: 2
%% type: 多选题
%% chapter: 电磁感应
%% point: 双杆问题
%% method: 无
%% score: 6
%% degree: 350
%% duplicateId: 0
%% body:
如图，两根光滑平行金属导轨固定在绝缘水平面上，左、右两侧导轨间距分别为 $d$ 和 $2d$，处于竖直向上的磁场中，磁感应强度大小分别为 $2B$ 和 $B$。已知导体棒 MN 的电阻为 $R$、长度为 $d$，导体棒 PQ 的电阻为 $2R$、长度为 $2d$，PQ 的质量是 MN 的 2 倍。初始时刻两棒静止，两棒中点之间连接一压缩量为 $L$ 的轻质绝缘弹簧。释放弹簧，两棒在各自磁场中运动直至停止，弹簧始终在弹性限度内。整个过程中两棒保持与导轨垂直并接触良好，导轨足够长且电阻不计。下列说法正确的是\xzanswer{AC}
\onepicture[w=7cm]{figs/10.png}
\fourchoices[answer=AC]
{弹簧伸展过程中、回路中产生顺时针方向的电流}
{PQ 速率为 $v$ 时，MN 所受安培力大小为 $\frac{{4{B^2}{d^2}v}}{{3R}}$}
{整个运动过程中，MN 与 PQ 的路程之比为 $2:1$}
{整个运动过程中，通过 MN 的电荷量为 $\frac{{BLd}}{{3R}}$}
%% answer: AC
%% memo:
\memoanswer{
	A．弹簧伸展过程中，根据右手定则可知，回路中产生顺时针方向的电流，选项 A 正确；

	B．任意时刻，设电流为 $I$，则 PQ 受安培力 $F_{PQ}=BI\cdot2d$

	方向向左；MN 受安培力 $F_{MN}=2BId$

	方向向右，可知两棒系统受合外力为零，动量守恒，设 PQ 质量为 $2m$，则 MN 质量为 $m$，PQ 速率为 $v$ 时，则 $2mv = mv'$

	解得 $v' = 2v$

	回路的感应电流 $I=\frac{2Bdv+B\cdot2dv}{3R}=\frac{2Bdv}{R}$

	MN 所受安培力大小为 $F_{MN}=2BId=\frac{4B^{2}d^{2}v}{R}$

	选项 B 错误；

	C．两棒最终停止时弹簧处于原长状态，由动量守恒可得

	$mx_{1}=2mx_{2}$

	$x_{1}+x_{2}=L$

	可得最终 MN 位置向左移动

	$x_{1}=\frac{2L}{3}$

	PQ 位置向右移动

	$x_{2}=\frac{L}{3}$

	因任意时刻两棒受安培力和弹簧弹力大小都相同，设整个过程两棒受的弹力的平均值为 $F_{\text{弹}}$，安培力平均值 $F_{\text{安}}$，则整个过程根据动能定理

	$F_{\text{弹}}x_{1}-F_{\text{安}}x_{MN}=0$

	$F_{\text{弹}}x_{2}-F_{\text{安}}x_{PQ}=0$

	可得

	$\frac{x_{MN}}{x_{PQ}}=\frac{x_{1}}{x_{2}}=\frac{2}{1}$

	选项 C 正确；

	D．两棒最后停止时，弹簧处于原长位置，此时两棒间距增加了 $L$，由上述分析可知，MN 向左位置移动 $\frac{2L}{3}$，PQ 位置向右移动 $\frac{L}{3}$，则

	$q=\overline{I}\Delta t=\frac{\Delta\Phi}{R_{\text{总}}}=\frac{2B\frac{2L}{3}d+B\frac{L}{3}2d}{3R}=\frac{2BLd}{3R}$

	选项 D 错误。

	故选 AC。
}

\item
%% number: 11
%% paperName: 2023年辽宁高考第 11 题 8 分
%% typeId: 4
%% type: 实验题
%% chapter: 运动和力
%% point: 验证动量守恒实验
%% method: 无
%% score: 8
%% degree: 550
%% duplicateId: 0
%% body:
某同学为了验证对心碰撞过程中的动量守恒定律，设计了如下实验：用纸板搭建如图所示的滑道，使硬币可以平滑地从斜面滑到水平面上，其中 OA 为水平段。选择相同材质的一元硬币和一角硬币进行实验。

测量硬币的质量，得到一元和一角硬币的质量分别为 $m_{1}$ 和 $m_{2}$（$m_{1} > m_{2}$）。将硬币甲放置在斜面一某一位置，标记此位置为 B。由静止释放甲，当甲停在水平面上某处时，测量甲从 O 点到停止处的滑行距离 OP。将硬币乙放置在 O 处，左侧与 O 点重合，将甲放置于 B 点由静止释放。当两枚硬币发生碰撞后，分别测量甲乙从 O 点到停止处的滑行距离 OM 和 ON。保持释放位置不变，重复实验若干次，得到 OP、OM、ON 的平均值分别为 $s_{0}$、$s_{1}$、$s_{2}$。
\begin{enumerate}
	\item
	在本实验中，甲选用的是\_\_\_\_（填“一元”或“一角”）硬币；
	\item
	碰撞前，甲到 O 点时速度的大小可表示为\_\_\_\_\_（设硬币与纸板间的动摩擦因数为 $\mu$，重力加速度为 $g$）；
	\item
	若甲、乙碰撞过程中动量守恒，则 $\frac{{\sqrt {{s_0}}  - \sqrt {{s_1}} }}{{\sqrt {{s_2}} }}$ = \_\_\_\_\_（用 $m_{1}$ 和 $m_{2}$ 表示），然后通过测得的具体数据验证硬币对心碰撞过程中动量是否守恒；
	\item
	由于存在某种系统或偶然误差，计算得到碰撞前后甲动量变化量大小与乙动量变化量大小的比值不是 1，写出一条产生这种误差可能的原因\_\_\_\_\_\_\_\_\_\_。
\end{enumerate}
\onepicture[w=14cm]{figs/11.png}
%% answer: 
\jdanswer{
\begin{enumerate}
	\item
	一元
	\item
	$\sqrt {2\mu g{s_0}}$
	\item
	$\frac{{{m_2}}}{{{m_1}}}$
	\item
	由于存在某种系统或偶然误差，计算得到碰撞前后甲动量变化量大小与乙动量变化量大小的比值不是 1，产生这种误差可能的原因有：测量误差，因为无论是再精良的仪器总是会有误差的，不可能做到绝对准确；碰撞过程中，我们认为内力远大于外力，动量守恒，实际上碰撞过程中，两个硬币组成的系统合外力不为零。
\end{enumerate}
}
%% memo:
\memoanswer{
\begin{enumerate}
	\item
	根据题意可知，甲与乙碰撞后没有反弹，可知甲的质量大于乙的质量，甲选用的是一元硬币。
	\item
	甲从 O 点到 P 点，根据动能定理 $-\mu m_{1}gs_{0}=0-\frac{1}{2}m_{1}v_{1}^{2}$，解得碰撞前，甲到 O 点时速度的大小 $v_{0} = \sqrt{2\mu gs_{0}}$。
	\item
	同理可得，碰撞后甲的速度和乙的速度分别为 $v_{1} = \sqrt{2\mu gs_{1}}$，$v_{2} = \sqrt{2\mu gs_{2}}$。若动量守恒，则满足 $m_{1}v_{0}=m_{1}v_{1}+m_{2}v_{2}$，整理可得 $\frac{\sqrt{s_{0}}-\sqrt{s_{1}}}{\sqrt{s_{2}}}=\frac{m_{2}}{m_{1}}$。
	\item
	由于存在某种系统或偶然误差，计算得到碰撞前后甲动量变化量大小与乙动量变化量大小的比值不是 1，产生这种误差可能的原因有：测量误差，因为无论是再精良的仪器总是会有误差的，不可能做到绝对准确；碰撞过程中，我们认为内力远大于外力，动量守恒，实际上碰撞过程中，两个硬币组成的系统合外力不为零。
\end{enumerate}
}

\item
%% number: 12
%% paperName: 2023年辽宁高考第 12 题 6 分
%% typeId: 4
%% type: 实验题
%% chapter: 电路
%% point: 测量电阻率
%% method: 无
%% score: 6
%% degree: 550
%% duplicateId: 0
%% body:
导电漆是将金属粉末添加于特定树脂原料中制作而成的能导电的喷涂油漆。现有一根用导电漆制成的截面为正方形的细长样品（固态），某同学欲测量其电阻率，设计了如图（a）所示的电路图，实验步骤如下：

a．测得样品截面的边长 $a=0.20\Ucm$；

b．将平行排列的四根金属探针甲、乙、丙、丁与样品接触，其中甲、乙、丁位置固定，丙可在乙、丁间左右移动；

c．将丙调节至某位置，测量丙和某探针之间的距离 $L$；

d．闭合开关 S，调节电阻箱 $R$ 的阻值，使电流表示数 $I=0.40\UA$，读出相应的电压表示数 $U$，断开开关 S；

e．改变丙的位置，重复步骤 c、d，测量多组 $L$ 和 $U$，作出 $U - L$ 图像如图（b）所示，得到直线的斜率 $k$。

回答下列问题：
\begin{enumerate}
	\item
	$L$ 是丙到\_\_\_\_\_\_\_\_\_\_\_（填“甲”“乙”或“丁”）的距离；
	\item
	写出电阻率的表达式 $\rho$ = \_\_\_\_\_\_\_\_\_\_\_（用 $k$、$a$、$I$ 表示）；
	\item
	根据图像计算出该样品的电阻率 $\rho$ = \_\_\_\_\_\_\_\_\_\_\_$\UOm$（保留两位有效数字）。
\end{enumerate}
\onepicture[w=13cm, align=right]{figs/12.png}
%% answer: 
\jdanswer{
\begin{enumerate}
	\item
	乙
	\item
	$\frac{{k{a^2}}}{I}$
	\item
	$6.5\times10^{-5}$
\end{enumerate}
}
%% memo:
\memoanswer{
\begin{enumerate}
	\item
	由于电压表测量的是乙、丙之间的电压，则 $L$ 是丙到乙的距离。
	\item
	根据电阻定律有 $R=\rho\frac{L}{a^{2}}$，再根据欧姆定律有 $R=\frac{U}{I}$，联立有 $U=\frac{\rho I}{a^{2}}L$，则 $\rho=\frac{ka^{2}}{I}$。
	\item
	根据图像可知 $k = 6.5\UVm$，则根据（2）代入数据有 $\rho = 6.5 \times 10^{-5}\UOm$。
\end{enumerate}
}

\item
%% number: 13
%% paperName: 2023年辽宁高考第 13 题 10 分
%% typeId: 6
%% type: 计算题
%% chapter: 直线运动
%% point: 匀变速直线运动
%% method: 无
%% score: 10
%% degree: 750
%% duplicateId: 0
%% body:
某大型水陆两栖飞机具有水面滑行汲水和空中投水等功能。某次演练中，该飞机在水面上由静止开始匀加速直线滑行并汲水，速度达到 $v_{1}=80\Ums$ 时离开水面，该过程滑行距离 $L=1\,600\Um$、汲水质量 $m=1.0\times10^{4}\Ukg$。离开水面后，飞机攀升高度 $h=100\Um$ 时速度达到 $v_{2}=100\Ums$，之后保持水平匀速飞行，待接近目标时开始空中投水。取重力加速度 $g=10\Umsq$。求：
\begin{enumerate}
	\item
	飞机在水面滑行阶段的加速度 $a$ 的大小及滑行时间 $t$；
	\item
	整个攀升阶段，飞机汲取的水的机械能增加量 $\Delta E$。
\end{enumerate}
%% answer: 
\jdanswer{
\begin{enumerate}
	\item
	$2\Umsq$，$40\Us$
	\item
	$6\times10^{7}\UJ$
\end{enumerate}
}
%% memo:
\memoanswer{
\begin{enumerate}
	\item
	飞机从静止开始做匀加速直线运动，平均速度为 $\frac{v_{1}}{2}$，则 $L = \frac{1}{2}v_{1}t$，解得飞机滑行的时间为 $t=\frac{2L}{v_{1}}=\frac{2\times1600}{80}\Us=40\Us$，飞机滑行的加速度为 $a=\frac{v_{1}}{t}=\frac{80}{40}\Umsq=2\Umsq$。
	\item
	飞机从水面至 $h = 100\Um$ 处，水的机械能包含水的动能和重力势能，则 $\Delta E = \frac{1}{2}mv_{2}^{2} + mgh = \frac{1}{2} \times 1 \times 10^{4} \times 100^{2}\UJ + 1 \times 10^{4} \times 10 \times 100\UJ = 6 \times 10^{7}\UJ$。
\end{enumerate}
}

\item
%% number: 14
%% paperName: 2023年辽宁高考第 14 题 13 分
%% typeId: 6
%% type: 计算题
%% chapter: 磁场
%% point: 带电粒子在组合场中的运动
%% method: 无
%% score: 13
%% degree: 450
%% duplicateId: 0
%% body:
如图，水平放置的两平行金属板间存在匀强电场，板长是板间距离的 $\sqrt 3 $ 倍。金属板外有一圆心为 O 的圆形区域，其内部存在磁感应强度大小为 $B$、方向垂直于纸面向外的匀强磁场。质量为 $m$、电荷量为 $q$（$q$ > 0）的粒子沿中线以速度 $v_{0}$ 水平向右射入两板间，恰好从下板边缘 P 点飞出电场，并沿 PO 方向从图中 O$'$ 点射入磁场。已知圆形磁场区域半径为 $\frac{{2m{v_0}}}{{3qB}}$，不计粒子重力。
\begin{enumerate}
	\item
	求金属板间电势差 $U$；
	\item
	求粒子射出磁场时与射入磁场时运动方向间的夹角 $\theta$；
	\item
	仅改变圆形磁场区域的位置，使粒子仍从图中 O$'$ 点射入磁场，且在磁场中的运动时间最长。定性画出粒子在磁场中的运动轨迹及相应的弦，标出改变后的侧形磁场区域的圆心 M。
\end{enumerate}
\onepicture[w=5.5cm, align=right]{figs/14.png}
%% answer: 
\jdanswer{
\begin{enumerate}
	\item
	$U = \frac{{mv_0^2}}{{3q}}$
	\item
	$\frac{\pi }{3}$ 或 $60^\circ$
	\item
	粒子在磁场中的运动轨迹及相应的弦、改变后磁场区域圆心 M 的位置如图所示。
\end{enumerate}
}
%% memo:
\memoanswer{
\begin{enumerate}
	\item
	设板间距离为 $d$，则板长为 $\sqrt{3}d$，带电粒子在板间做类平抛运动，两板间的电场强度为 $E = \frac{U}{d}$，根据牛顿第二定律得，电场力提供加速度 $qE = ma$，解得 $a = \frac{qU}{md}$。设粒子在平板间的运动时间为 $t_{0}$，根据类平抛运动的运动规律得 $\frac{d}{2} = \frac{1}{2}at_{0}^{2}$，$\sqrt{3}d = v_{0}t_{0}$，联立解得 $U = \frac{mv_{0}^{2}}{3q}$。
	\item
	设粒子出电场时与水平方向夹角为 $\alpha$，则有 $\tan\alpha=\frac{at_{0}}{v_{0}}=\frac{\sqrt{3}}{3}$，故 $\alpha=\frac{\pi}{6}$。则出电场时粒子的速度为 $v=\frac{v_{0}}{\cos\alpha}=\frac{2\sqrt{3}}{3}v_{0}$。粒子出电场后沿直线匀速直线运动，接着进入磁场，根据牛顿第二定律，洛伦兹力提供匀速圆周运动所需的向心力得 $qvB=m\frac{v^{2}}{r}$，$r=\frac{mv}{qB}=\frac{2\sqrt{3}mv_{0}}{3qB}$。已知圆形磁场区域半径为 $R = \frac{2mv_{0}}{3qB}$，故 $r = \sqrt{3}R$。粒子沿 PO 方向射入磁场即沿半径方向射入磁场，故粒子将沿半径方向射出磁场，粒子射出磁场时与射入磁场时运动方向的夹角为 $\theta$，则粒子在磁场中运动圆弧轨迹对应的圆心角也为 $\theta$，由几何关系可得 $\theta = 2\alpha = \frac{\pi}{3}$。故粒子射出磁场时与射入磁场时运动方向的夹角为 $\frac{\pi}{3}$ 或 $60^\circ$；

	\onepicture[h=4cm]{figs/14_an1.jpg}

	\item
	带电粒子在该磁场中运动的半径与圆形磁场半径关系为 $r = \sqrt{3}R$，根据几何关系可知，带电粒子在该磁场中运动的轨迹一定为劣弧，故劣弧所对应轨迹圆的弦为磁场圆的直径时粒子在磁场中运动的时间最长。则相对应的运动轨迹和弦以及圆心 M 的位置如图所示：

	\onepicture[h=4cm]{figs/14_an2.png}
\end{enumerate}
}

\item
%% number: 15
%% paperName: 2023年辽宁高考第 15 题 17 分
%% typeId: 6
%% type: 计算题
%% chapter: 运动和力
%% point: 动量和能量
%% method: 无
%% score: 17
%% degree: 350
%% duplicateId: 0
%% body:
如图，质量 $m_{1}=1\Ukg$ 的木板静止在光滑水平地面上，右侧的竖直墙面固定一劲度系数 $k=20\UNm$ 的轻弹簧，弹簧处于自然状态。质量 $m_{2}=4\Ukg$ 的小物块以水平向右的速度 $v_{0}=\frac{5}{4}\Ums$ 滑上木板左端，两者共速时木板恰好与弹簧接触。木板足够长，物块与木板间的动摩擦因数 $\mu=0.1$，最大静摩擦力等于滑动摩擦力。弹簧始终处在弹性限度内，弹簧的弹性势能 $E_{p}$ 与形变量 $x$ 的关系为 $E_{p}=\frac{1}{2}kx^{2}$。取重力加速度 $g=10\Umsq$，结果可用根式表示。
\begin{enumerate}
	\item
	求木板刚接触弹簧时速度 $v_{1}$ 的大小及木板运动前右端距弹簧左端的距离 $x_{1}$；
	\item
	求木板与弹簧接触以后，物块与木板之间即将相对滑动时弹簧的压缩量 $x_{2}$ 及此时木板速度 $v_{2}$ 的大小；
	\item
	已知木板向右运动的速度从 $v_{2}$ 减小到 0 所用时间为 $t_{0}$。求木板从速度为 $v_{2}$ 时到之后与物块加速度首次相同时的过程中，系统因摩擦转化的内能 $\Delta U$（用 $t$ 表示）。
\end{enumerate}
\onepicture[w=13cm, align=right]{figs/15.png}
%% answer: 
\jdanswer{
\begin{enumerate}
	\item
	$1\Ums$；$0.125\Um$
	\item
	$0.25\Um$；$\frac{\sqrt{3}}{2}\Ums$
	\item
	$4\sqrt{3}t_{0}-8t_{0}^{2}$
\end{enumerate}
}
%% memo:
\memoanswer{
\begin{enumerate}
	\item
	由于地面光滑，则 $m_{1}$、$m_{2}$ 组成的系统动量守恒，则有 $m_{2}v_{0}=(m_{1}+m_{2})v_{1}$，代入数据有 $v_{1}=1\Ums$。对 $m_{1}$ 受力分析有 $a_{1}=\frac{\mu m_{2}g}{m_{1}}=4\Umsq$，则木板运动前右端距弹簧左端的距离有 $v_{1}^{2}=2a_{1}x_{1}$，代入数据解得 $x_{1}=0.125\Um$。
	\item
	木板与弹簧接触以后，对 $m_{1}$、$m_{2}$ 组成的系统有 $kx=(m_{1}+m_{2})a_{\text{共}}$，对 $m_{2}$ 有 $a_{2}=\mu g=1\Umsq$。当 $a_{\text{共}}=a_{2}$ 时物块与木板之间即将相对滑动，解得此时的弹簧压缩量 $x_{2}=0.25\Um$。对 $m_{1}$、$m_{2}$ 组成的系统列动能定理有 $-\frac{1}{2}kx_{2}^{2}=\frac{1}{2}(m_{1}+m_{2})v_{2}^{2}-\frac{1}{2}(m_{1}+m_{2})v_{1}^{2}$，代入数据有 $v_{2}=\frac{\sqrt{3}}{2}\Ums$。
	\item
	木板从速度为 $v_{2}$ 时到之后与物块加速度首次相同时的过程中，由于木板即 $m_{1}$ 的加速度大于木块 $m_{2}$ 的加速度，则当木板与木块的加速度相同时即弹簧形变量为 $x_{2}$ 时，则说明此时 $m_{1}$ 的速度大小为 $v_{2}$，共用时 $2t_{0}$，且 $m_{2}$ 一直受滑动摩擦力作用，则对 $m_{2}$ 有 $-\mu m_{2}g\cdot2t_{0}=m_{2}v_{3}-m_{2}v_{2}$，解得 $v_{3}=\frac{\sqrt{3}}{2}-2t_{0}$。则对于 $m_{1}$、$m_{2}$ 组成的系统有 $-W_{f}=\frac{1}{2}m_{1}v_{2}^{2}+\frac{1}{2}m_{2}v_{3}^{2}-\frac{1}{2}(m_{1}+m_{2})v_{2}^{2}$，$\Delta U = -W_{f}$，联立有 $\Delta U = 4\sqrt{3}t_{0}-8t_{0}^{2}$。
\end{enumerate}
}

\end{enumerate}

\end{document}
